\chapter{Disorientation}

\section*{Definition}
Impairment of awareness of time, place, or person, often accompanied by confusion, difficulty concentrating, and inability to perform familiar tasks.

\section*{Pathophysiology}
Disorientation results from dysfunction in cortical and subcortical networks supporting orientation and attention. Acute disorientation (delirium) involves global cognitive impairment from metabolic, toxic, infectious, or structural causes. Chronic disorientation (dementia) reflects progressive neurodegenerative damage. The reticular activating system and prefrontal cortex are key affected regions.

\section*{Causes}
\begin{itemize}
  \item Delirium (metabolic, infectious, medication-induced)
  \item Alzheimer disease and other dementias
  \item Head trauma or traumatic brain injury
  \item Stroke or transient ischemic attack
  \item Seizure (especially complex partial or post-ictal)
  \item Hypoglycemia or hyperglycemia
  \item Electrolyte imbalances
  \item Medication side effects (anticholinergics, benzodiazepines, opioids)
  \item Alcohol intoxication or withdrawal
  \item Severe dehydration or malnutrition
\end{itemize}

\section*{When to See a Doctor}
See your doctor if:
\begin{itemize}
  \item Severe or worsening symptoms
  \item Sudden-onset disorientation (possible stroke, hypoglycemia, or delirium — emergency evaluation)
  \item Accompanied by other concerning signs
\end{itemize}

\section*{Related Symptoms}
\begin{itemize}
  \item Memory loss
  \item Confusion
  \item Agitation
  \item Dizziness
  \item Speech difficulty
\end{itemize}
\textbf{Important:} If this symptom persists, your doctor will check for serious underlying causes.

\section*{Self-Care / Home Management}
\begin{itemize}
  \item Rest and avoid activities that make it worse.
  \item Maintain adequate hydration and a balanced diet.
  \item Over-the-counter remedies may provide symptomatic relief (consult a pharmacist or doctor).
  \item Monitor symptoms and seek professional advice if they persist or worsen.
  \item Avoid self-medicating with prescription medications without medical guidance.
\end{itemize}

\section*{How Your Doctor Will Evaluate You}
\begin{itemize}
  \item A thorough history and physical examination are the first steps in evaluation.
  \item Laboratory tests (blood work, urinalysis) may be ordered based on clinical suspicion.
  \item Imaging studies (X-ray, ultrasound, CT, MRI) may be indicated when structural pathology is suspected.
  \item Specialist referral is arranged when the diagnosis remains uncertain or requires specific expertise.
\end{itemize}

\section*{Treatment Options}
\begin{itemize}
  \item Treatment targets the underlying cause identified through diagnostic evaluation.
  \item Supportive care and symptom management are provided while awaiting definitive therapy.
  \item Medications, lifestyle modifications, and physical therapies may be prescribed as appropriate.
  \item Follow-up ensures treatment effectiveness and allows timely adjustments to the care plan.
\end{itemize}

\clearpage